\section{Struktur Organisasi}

Struktur organisasi di PT Inti Utama Solusindo (IUS) dirancang untuk mendukung
proses pengembangan solusi teknologi secara terintegrasi bagi ekosistem bisnis
Pharos Group. Setiap divisi memiliki peran dan tanggung jawab yang spesifik,
mulai dari perumusan strategi teknologi, pengembangan software, hingga dukungan
operasional sistem yang telah diimplementasikan. Pembagian struktur ini
memungkinkan proses software development berjalan secara terorganisir,
kolaboratif, serta mendukung keberlanjutan pengembangan produk digital
perusahaan.

\begin{figure}[H]
      \centering
      \includegraphics[width=14cm]{struktur-organisasi.png}
      \caption{Struktur Organisasi PT Inti Utama Solusindo}%
\end{figure}

\subsection{Direksi}
Posisi pimpinan tertinggi di PT Inti Utama Solusindo dipimpin oleh Brian
Lembong selaku Direksi. Direksi bertanggung jawab dalam menetapkan arah
strategis perusahaan, khususnya dalam pengembangan teknologi yang mendukung
berbagai unit bisnis dalam Pharos Group.

Selain berperan dalam pengambilan keputusan strategis, Direksi juga memiliki
keterlibatan dalam pengembangan sumber daya manusia di bidang teknologi. Salah
satu bentuk keterlibatan tersebut adalah partisipasi langsung dalam proses
supervisi dan evaluasi talenta baru, termasuk dalam tahapan seleksi teknis
mahasiswa magang yang akan bergabung dengan tim pengembang di perusahaan.

\subsection{\textit{Developer}}
Divisi utama yang bertanggung jawab terhadap pengembangan produk digital di PT Inti Utama Solusindo.
Divisi ini berada di bawah kepemimpinan Michael Tjitra selaku \textit{Manager}
dan berperan dalam mengelola seluruh proses pengembangan aplikasi, mulai dari implementasi
fitur, pengujian sistem, hingga pemeliharaan \textit{codebase}.

Struktur internal divisi ini terdiri dari beberapa tim dengan fungsi yang
berbeda, yaitu:

{\leftskip=1.27cm
\subsubsection{\textbf{a. Tim Software Development}}
Tim Software Development terbagi menjadi enam kelompok pengembang yang
masing-masing dipimpin oleh seorang Supervisor. Beberapa supervisor yang
memimpin tim pengembang antara lain Terry Sutartho, M. Fariz Taswarul Afkar,
Muhammad Ilham Egi, Dimas Aditya Nugroho, dan Kelvin Samuel.

Pembagian tim ini memungkinkan perusahaan untuk mengelola berbagai modul dan
fitur aplikasi secara paralel. Setiap tim bertanggung jawab atas pengembangan
bagian tertentu dari sistem, sehingga proses development dapat berjalan lebih
efisien serta memudahkan koordinasi dalam pengelolaan source code dan
implementasi fitur baru.

\subsubsection{\textbf{b. Tim Quality Assurance (QA)}}

Tim Quality Assurance (QA) dipimpin oleh Rachelle Marvella dan memiliki
tanggung jawab dalam melakukan proses pengujian terhadap aplikasi yang
dikembangkan oleh tim developer. Pengujian ini mencakup berbagai aspek seperti
fungsionalitas sistem, stabilitas aplikasi, serta kesesuaian implementasi fitur
dengan requirement yang telah ditentukan.

Melalui proses testing yang sistematis, tim QA memastikan bahwa setiap
perubahan pada code telah melalui proses validasi sebelum dirilis ke tahap
produksi. Hal ini penting untuk meminimalkan potensi bug serta menjaga kualitas
sistem yang digunakan oleh pengguna.

\subsubsection{\textbf{c. Tim Artificial Intelligence Development}}

Selain tim pengembangan utama, divisi ini juga memiliki Tim Artificial
Intelligence Development yang dipimpin oleh Devany Aguslia (Supervisor-Intern).
Tim ini berfokus pada eksplorasi dan pengembangan teknologi kecerdasan buatan
yang berpotensi diintegrasikan ke dalam sistem digital perusahaan di masa
mendatang.

Inisiatif ini menunjukkan komitmen perusahaan dalam mengikuti perkembangan
teknologi terbaru, khususnya dalam pengembangan aplikasi kesehatan yang semakin
memanfaatkan kemampuan analisis data dan otomatisasi berbasis AI.

}

\subsection{\textit{Technical Support}}

Struktur organisasi di PT Inti Utama Solusindo menunjukkan penerapan pola kerja
yang terorganisir dengan pembagian peran yang jelas pada setiap fungsi
teknologi. Setiap divisi memiliki tanggung jawab yang saling melengkapi, mulai
dari tahap pengembangan software, proses pengujian kualitas sistem, hingga
dukungan operasional setelah sistem diimplementasikan.
